\subsection{Analisis Waktu--Frekuensi Menggunakan STFT}
FFT global belum membedakan apakah puncak kuat berasal dari nada
panjang atau banyak kejadian pendek. STFT menempatkan komponen
tersebut pada sumbu waktu. Digunakan Hann 1024 dan 4096 sampel,
hop seperempat panjang jendela, serta hanya bingkai lengkap tanpa
penambahan nol pada ujung sinyal.

Gambar~\ref{fig:stftbass}--\ref{fig:stfttuba} menggunakan interval
29--33 s, frekuensi 0--2000 Hz, dan warna $-110$ sampai $-35$ dB
dengan acuan sama. X adalah pusat jendela, bukan awal potongan.
Y adalah frekuensi; warna lebih terang berarti daya per Hz lebih
besar. Penskalaan terhadap energi jendela dijelaskan pada lampiran.
Spektrogram seluruh rekaman tersedia dalam notebook. Pembesaran ini
tidak diagregasi sehingga sesuai untuk membandingkan detail.

\clearpage
\subsubsection{Karakteristik Waktu--Frekuensi Bass}
Untuk melihat apakah energi rendah bertahan atau berubah menurut
not, digunakan spektrogram bass pada Gambar~\ref{fig:stftbass}.
\begin{figure}[H]
\centering\includegraphics[width=\linewidth]{stft_zoom_Bass.png}
\caption{Bass, 29--33 s. Atas: 1024 sampel; bawah: 4096 sampel.
Garis rendah yang berganti dan guratan singkat menunjukkan perubahan
kejadian pada bagian yang sama.}\label{fig:stftbass}
\end{figure}
Daerah terang terkonsentrasi pada frekuensi rendah, sesuai FFT.
Sekitar 29,25 s dan kejadian berikutnya, tampak guratan menuju
frekuensi lebih tinggi yang kemudian melemah. Pola ini sesuai
perubahan amplitudo pada awal not, bukan nada tinggi baru yang
harus dianggap bertahan. Panel 4096 memisahkan beberapa garis di
bawah 250 Hz; pada panel 1024 daerah itu lebih menyatu.
Secara matematis, perkalian sinyal dengan jendela menyebabkan
spektrumnya melebar menurut respons jendela. Ketebalan pita gambar
bukan ukuran jumlah nada sebenarnya.

\clearpage
\subsubsection{Karakteristik Waktu--Frekuensi Piano}
Pada piano, pertanyaannya adalah bagaimana rangkuman banyak not
dalam FFT berubah saat diamati lokal. STFT memperlihatkan lapisan
yang bertahan sekaligus pergantian kejadian pendek.
\begin{figure}[H]
\centering\includegraphics[width=\linewidth]{stft_zoom_Piano.png}
\caption{Piano, 29--33 s. Beberapa komponen tumpang tindih karena
kelompok piano berasal dari tiga channel sumber.}\label{fig:stftpiano}
\end{figure}
Sekitar 31,6--32,3 s, jalur terang bergerak naik dari beberapa ratus
Hz, sementara garis lain masih bertahan. Komposisi frekuensi berubah,
bukan seluruh spektrum sekadar bertambah amplitudonya. Panel 4096
menampilkan garis lebih tipis, sedangkan panel 1024 mengamati
perubahan waktu dalam potongan lebih pendek. Karena not bersamaan,
satu garis dapat merupakan fundamental suatu not atau harmonik
not lain. Linearitas Fourier menjelaskan penjumlahan kontribusi,
tetapi grafik daya saja tidak memisahkannya secara unik.

\clearpage
\subsubsection{Karakteristik Waktu--Frekuensi Gitar}
Kelompok gitar menggabungkan pattern Guitar dan Guitar Melody.
Analisis lokal diperlukan agar perubahan aransemen dibedakan dari
sifat preset gitar itu sendiri.
\begin{figure}[H]
\centering\includegraphics[width=\linewidth]{stft_zoom_Guitar.png}
\caption{Gitar, 29--33 s. Celah gelap berarti daya lebih rendah menurut
batas warna, bukan selalu keheningan mutlak.}\label{fig:stftgitar}
\end{figure}
Sekitar 31,6--32,3 s, jalur terang di frekuensi menengah bergerak naik.
Setelahnya tampak selang berdaya lebih kecil antara kejadian berikut.
Panel pendek memperjelas perubahan antarbagian; panel panjang
merinci garis dalam suatu kejadian. Melodi dan iringan berkontribusi
pada spektrum kelompok yang sama. Secara matematis, perubahan
osilasi dan envelope mengubah posisi serta penyebaran spektrum
lokal. Pola ini tidak membuktikan jenis gitar berubah selama lagu.

\clearpage
\subsubsection{Karakteristik Waktu--Frekuensi Brass}
Brass memberi contoh berbeda karena akor kedua berlangsung dari
29,23074 sampai 32,30766 s. STFT digunakan untuk memeriksa apakah
komponen frekuensinya bertahan selama akor tersebut.
\begin{figure}[H]
\centering\includegraphics[width=\linewidth]{stft_zoom_Brass.png}
\caption{Brass pada akor kedua. Struktur mendatar menunjukkan
komponen yang bertahan; warna melemah setelah note-off.}\label{fig:stftbrass}
\end{figure}
Garis terang sekitar 0,59 dan 1,17 kHz bertahan pada sebagian besar
interval akor. Pada 1,4--1,6 kHz, panel 4096 memperlihatkan beberapa
garis terpisah, sedangkan panel 1024 lebih menyatu. Ini bukti visual
spesifik pengaruh panjang jendela pada pemisahan frekuensi.
Warna di awal dan akhir tidak berhenti tepat pada waktu not karena
jendela melibatkan sampel di sekitar pusatnya, serta bunyi memiliki
attack dan release. Pola mendatar sesuai not yang ditahan, tetapi
bukan bukti amplitudo benar-benar konstan. Osilasi yang bertahan
memusatkan frekuensi; pembatasan waktunya menimbulkan pelebaran.

\clearpage
\subsubsection{Karakteristik Waktu--Frekuensi Tuba}
Bass dan tuba menarik dibandingkan karena seluruh onset not sumber
sama, sedangkan 20 dari 187 note-off berbeda 2 atau 4 tick.
STFT tuba memperlihatkan pengaruh pitch satu oktaf lebih tinggi dan
preset berbeda terhadap tampilan lokal.
\begin{figure}[H]
\centering\includegraphics[width=\linewidth]{stft_zoom_Tuba.png}
\caption{Tuba, 29--33 s. Kejadian mengikuti onset not yang sama dengan
bass, tetapi distribusi frekuensinya tidak identik.}\label{fig:stfttuba}
\end{figure}
Daya tuba tersebar pada beberapa garis rendah hingga menengah,
sedangkan bass lebih terkonsentrasi di bawah 250 Hz. Pergantian
kejadian mengikuti onset yang sama, tetapi warna dan kekuatan
harmonik berbeda. Hubungan oktaf menjelaskan frekuensi nada ideal;
pembobotan harmonik menjelaskan mengapa seluruh spektrum bukan
sekadar salinan. Preset GM Tuba tidak mereproduksi artikulasi FLEX
\textit{Tubas Staccato} secara identik, sehingga interpretasi dibatasi
pada audio rekonstruksi.

\subsection{Pengaruh Panjang Jendela terhadap Resolusi Waktu dan Frekuensi}
Perbandingan ini ingin mengetahui informasi yang bertambah atau
berkurang ketika potongan diperpanjang. Panjang 1024 dan 4096 dipilih
karena berbeda empat kali, sementara bentuk Hann dan proporsi
tumpang tindih sama. Perubahan dapat dikaitkan dengan panjang jendela.
\begin{table}[H]
\centering\small
\caption{Parameter pada $f_s=44\,100$ Hz; overlap keduanya 75\%.}
\label{tab:jendela}
\begin{tabular}{lrrrrr}
\toprule
$M$ & $H$ & Durasi (ms) & Hop (ms) & $\Delta\nu$ (Hz) & ENBW (Hz) \\
\midrule
1024 & 256 & 23.220 & 5.805 & 43.066 & 64.600 \\
4096 & 1024 & 92.880 & 23.220 & 10.767 & 16.150 \\
\bottomrule
\end{tabular}

\end{table}
Durasi bertambah dari 23,22 menjadi 92,88 ms, sedangkan jarak bin
turun dari 43,07 menjadi 10,77 Hz. Lebar lobus utama Hann antara
dua nol pertama sekitar 172,27 dan 43,07 Hz. Ini menunjukkan
mengapa jarak bin bukan satu-satunya ukuran resolusi.
ENBW adalah lebar pita ekuivalen energi respons jendela, dihitung
dengan $f_s\sum w_r^2/(\sum w_r)^2=1,5f_s/M$.

Pada brass, garis sekitar 1,4--1,6 kHz terpisah lebih jelas untuk
4096. Sebaliknya, pergantian kejadian piano dan gitar melibatkan
rentang sampel lebih luas. Ambang daya 10\% di sekitar awal akor
kedua brass tercapai pada pusat bingkai sekitar 29,24553 s untuk
1024 dan 29,23392 s untuk 4096. Jendela panjang yang tampak lebih
awal tidak berarti onset sebenarnya lebih awal: bingkai terpusat
sudah mencakup sampel setelah pusat. Angka ambang ini bukan ukuran
galat pendeteksi onset terhadap suatu standar uji.

Memperpanjang jendela empat kali mempersempit respons frekuensi,
tetapi memperlebar dukungan waktu. Mengurangi hop saja memperpadat
kisi waktu tanpa memperpendek jendela. Penambahan nol pada FFT juga
hanya memperpadat kisi frekuensi tanpa mempersempit lobus utama.
Model Hann kontinu menghasilkan
$\sigma_t\sigma_\omega\approx0,513117>1/2$. Nilai ini menjelaskan
penyebaran energi jendela, bukan ketelitian aktual pendeteksi drum
atau not musik.

\subsection{Perubahan Temporal dan Karakteristik Nonstasioner}
Analisis ini membedakan perubahan amplitudo dari perubahan komposisi
frekuensi. Rekaman dibagi menjadi segmen 10 s; RMS, sentroid dan
fraksi pita dihitung dari periodogram setiap segmen. Segmen hampir
senyap tidak ditafsirkan melalui sentroidnya. Tabel berikut memilih
contoh yang mempunyai amplitudo bermakna.
\begin{table}[H]
\centering\small
\caption{Contoh perubahan temporal pada WAV baru.}\label{tab:temporal}
\begin{tabular}{llrrr}
\toprule
Instrumen & Interval (s) & RMS & Sentroid (Hz) & $\geq1000$ Hz (\%) \\
\midrule
Piano & 0--10 & 0.01073 & 692.76 & 10.06 \\
Piano & 50--60 & 0.01744 & 536.02 & 6.49 \\
Guitar & 0--10 & 0.00561 & 376.61 & 6.21 \\
Guitar & 30--40 & 0.01330 & 765.47 & 15.77 \\
Bass & 10--20 & 0.00950 & 90.43 & 0.04 \\
Bass & 30--40 & 0.01055 & 91.91 & 0.04 \\
Tuba & 10--20 & 0.01044 & 304.16 & 0.52 \\
Tuba & 30--40 & 0.01138 & 313.70 & 0.43 \\
\bottomrule
\end{tabular}
\end{table}
Piano pada 0--10 s mempunyai RMS lebih kecil daripada 50--60 s,
tetapi sentroid dan proporsi frekuensi tingginya lebih besar.
Bagian beramplitudo efektif kecil dapat memiliki komposisi spektral
relatif lebih tinggi. Keras-pelan secara amplitudo berbeda dari
tinggi-rendahnya distribusi frekuensi.

Pada gitar, sentroid naik dari sekitar 376,61 menjadi 765,47 Hz
antara 0--10 dan 30--40 s. Pusat daya bergeser ke frekuensi tinggi,
sesuai peran melodi dan register pada data serta spektrogram.
Fraksi di atas 1000 Hz naik dari 6,21\% menjadi 15,77\%.
Pengalian seluruh sinyal dengan konstanta tidak mengubah fraksi
pita, sehingga hasil ini bukan sekadar perubahan volume seragam.

Bass mempunyai ringkasan pita lebih stabil daripada gitar, tetapi
not dan amplitudo lokalnya tetap berubah. Tuba juga berganti not
meskipun sentroid antarsegmen lebih dekat. Brass berganti dari
diam menjadi dua kelompok akor. Gabungan waveform, not dan
spektrogram mendukung deskripsi nonstasioner tiap rekaman: karakter
amplitudo atau frekuensinya berubah. Ini bukan uji formal proses
stokastik. Pada proses stasioner, sifat statistik menurut definisi
yang dipilih harus tetap terhadap pergeseran waktu; hal itu tidak
dibuktikan atau ditolak hanya dari tabel satu rekaman deterministik.

\subsection{Analisis Kanal DRUMS dan VOCALS dalam Format WAV}
Analisis tambahan ini membandingkan dua kanal audio yang sudah
tersedia, khususnya pola kejadian singkat dan komponen yang lebih
jarang muncul. Sumber kedua kanal berbeda dari lima instrumen
SoundFont, sehingga perolehan data dijelaskan sebelum hasilnya.

\subsubsection{Konversi Format dan Validasi Sampel}
Channel 11 dan 12 pada FLP masing-masing merujuk ke
\texttt{DRUMS.mp3} dan \texttt{VOCALS.mp3}. Tidak terdapat not pattern
untuk merender ulang kedua kanal dengan MIDI. Karena itu, audio
didekode lalu disimpan sebagai WAV float32 tanpa resampling,
normalisasi, atau perubahan amplitudo menggunakan SoundFile
\cite{soundfile}. Masing-masing berisi 7.090.031 frame stereo pada
44.100 Hz, berdurasi 160,771678 s.

MP3 merupakan sumber yang telah mengalami kompresi lossy. WAV hasil
konversi hanya menyimpan hasil dekodenya tanpa kompresi lossy baru;
detail yang hilang tidak dipulihkan. Pemeriksaan baca-tulis
menunjukkan kesamaan setiap sampel float32, dengan galat maksimum
nol. Kedua WAV tidak senyap dan tidak memiliki sampel dengan
magnitudo mencapai satu. Puncak amplitudo stereo DRUMS sekitar
0,89477, sedangkan VOCALS 0,99853.

Konsekuensi matematisnya sederhana: jika vektor sampel sama,
RMS, DFT dan STFT dengan parameter sama juga sama. Analisis DRUMS
dihitung ulang dari WAV; RMS, energi dan puncak frekuensinya sama
dengan perhitungan langsung dari dekode MP3. Dengan demikian, penggunaan WAV
di sini adalah pilihan format kerja, bukan peningkatan ketelitian
analisis atau kualitas sumber.

Parameter stretching yang terbaca pada kedua channel adalah
multiplier satu, pitch nol dan time nol, dengan start offset
parameter nol. Namun, pemetaan offset klip audio ke sampel serta
hasil playback belum diperiksa terhadap render FL Studio acuan.
Analisis memakai berkas penuh, tanpa memotong atau memaksanya menjadi
69,23 s. Kedua kanal tidak dimasukkan ke
\texttt{All\_Instruments.wav}. Nama VOCALS mengikuti label sumber
dan tidak dianggap sebagai bukti bahwa semua komponennya suara
manusia murni.

\subsubsection{Aktivitas Temporal dan Amplitudo Efektif}
Perbandingan waktu bertujuan mengetahui apakah kedua kanal aktif
sepanjang rekaman atau hanya pada beberapa bagian. Waveform dan
RMS dihitung dengan metode yang sama seperti instrumen rekonstruksi.
\begin{figure}[H]
\centering\includegraphics[width=\linewidth]{01_waveform_rms_kanal.png}
\caption{Amplitudo dan RMS kanal audio penuh. X: waktu dari awal
berkas sumber; kiri: amplitudo digital; kanan: RMS 100 ms.
Skala antarbaris sama.}\label{fig:kanalwaktu}
\end{figure}
Gambar~\ref{fig:kanalwaktu} memperlihatkan DRUMS aktif pada sebagian
besar rekaman, sedangkan VOCALS mempunyai kelompok kejadian utama
sekitar 31,4--35,2 s, 92,3--100,0 s dan 153,0--155,9 s.
Waktu ini mengikuti batas blok aktivitas, bukan transkripsi suku
kata atau waktu onset yang persis. Dengan ambang 1\% RMS blok
maksimum tiap kanal, durasi aktif DRUMS sekitar 156,1 s dan VOCALS
14,2 s.

RMS global VOCALS sekitar 0,04895, tetapi RMS pada blok aktif sekitar
0,16470. Untuk DRUMS, keduanya lebih dekat, yaitu 0,07757 dan 0,07872.
Bagian diam menurunkan rerata kuadrat sepanjang rekaman. Karena itu,
RMS global VOCALS yang lebih kecil tidak berarti setiap kejadiannya
beramplitudo lebih kecil daripada DRUMS. Perubahan dari latar
beramplitudo rendah menuju kejadian kuat juga mendukung deskripsi
nonstasioner, tanpa menyatakan hasil uji statistik formal.

\subsubsection{Spektrum Global dan Distribusi Energi}
FFT digunakan untuk membandingkan wilayah frekuensi yang paling
berkontribusi pada kedua kanal. Semua sampel berkas penuh digunakan;
jarak bin sekitar 0,006220 Hz. Grafik berikut memperlihatkan
0--5000 Hz, sedangkan perhitungan energi tetap sampai Nyquist.
\begin{figure}[H]
\centering\includegraphics[width=.96\linewidth]{02_fft_kanal.png}
\caption{Kerapatan daya DRUMS dan VOCALS dengan acuan sama.
X: frekuensi; Y: daya per Hz dalam dB. Tampilan merangkum maksimum
per 0,5 Hz; perhitungan memakai bin asli.}\label{fig:kanalfft}
\end{figure}
\begin{table}[H]
\centering\small
\caption{Hasil perhitungan dari dua WAV hasil dekode.}
\begin{tabular}{lrrrrr}
\toprule
Kanal & RMS & Puncak (Hz) & Sentroid (Hz) & $<250$ Hz (\%) & $\geq1000$ Hz (\%) \\
\midrule
DRUMS & 0.07757 & 41.60 & 563.50 & 91.19 & 6.83 \\
VOCALS & 0.04895 & 780.08 & 1195.31 & 0.02 & 31.87 \\
\bottomrule
\end{tabular}

\end{table}
DRUMS memiliki sekitar 91,19\% energi di bawah 250 Hz, sedangkan
VOCALS hanya sekitar 0,017\%. Pada VOCALS, sekitar 68,11\% berada
pada 250--1000 Hz dan 31,87\% pada 1000 Hz ke atas. Komposisi
VOCALS lebih banyak berkontribusi di wilayah menengah dan tinggi
dibandingkan DRUMS. Ini perbedaan distribusi spektral, bukan bukti
perbedaan keras suara perseptual atau kemurnian jenis sumber.

Puncak terpilih VOCALS sekitar 780,08 Hz tidak otomatis merupakan
frekuensi dasar suara manusia. Komponen yang bertahan, harmonik,
dan perubahan nada sama-sama berkontribusi pada FFT global.
Sentroid VOCALS sekitar 1195,31 Hz, dibandingkan DRUMS 563,50 Hz,
menunjukkan pusat distribusi daya pada frekuensi lebih tinggi.
Parseval kedua WAV menghasilkan galat relatif sekitar
$2,1\times10^{-15}$ atau lebih kecil. Hasil DRUMS lama tetap
tersimpan dan tidak digunakan sebagai pengganti perhitungan VOCALS.

\clearpage
\subsubsection{Karakteristik Waktu--Frekuensi DRUMS}
STFT dipakai untuk melokalisasi kejadian pada spektrum DRUMS.
Gambar~\ref{fig:stftdrumswav} dan~\ref{fig:stftvocalswav} menggunakan
29--33 s menurut waktu berkas sumber, frekuensi 0--5000 Hz, serta
skala warna yang sama, $-110$ sampai $-20$ dB. Kedua panjang Hann
tetap 1024 dan 4096, dengan overlap 75\%. Batas tampilan ini berbeda
dari panel lima instrumen; perbandingan warna langsung dilakukan
di dalam pasangan yang memakai skala sama.
\begin{figure}[H]
\centering\includegraphics[width=\linewidth]{stft_zoom_DRUMS.png}
\caption{STFT DRUMS dari WAV hasil dekode, pada waktu sumber 29--33 s.
Atas: 1024 sampel; bawah: 4096 sampel.}\label{fig:stftdrumswav}
\end{figure}
X adalah pusat jendela; Y adalah frekuensi. Daerah terang rendah
disertai guratan yang menjangkau frekuensi lebih tinggi, terutama
pada kejadian berulang sekitar 29--31 s. Pola itu sesuai dengan
perubahan amplitudo singkat yang melibatkan banyak frekuensi,
bukan hanya satu osilasi tetap. Panel 4096 memperinci garis
horizontal, sementara panel 1024 melibatkan potongan waktu lebih
pendek. Perkalian dengan Hann menjelaskan kompromi antara
penyebaran frekuensi dan dukungan waktu. Grafik ini tidak dipakai
untuk mengidentifikasi jenis setiap pukulan drum secara otomatis.

\clearpage
\subsubsection{Karakteristik Waktu--Frekuensi VOCALS}
STFT VOCALS memperlihatkan hubungan antara periode beramplitudo
rendah dan kejadian yang membentuk puncak pada FFT global.
\begin{figure}[H]
\centering\includegraphics[width=\linewidth]{stft_zoom_VOCALS.png}
\caption{STFT VOCALS pada waktu sumber 29--33 s. Komponen kuat mulai
tampak sekitar 31,4 s. Nama kanal mengikuti sumber FLP.}
\label{fig:stftvocalswav}
\end{figure}
Sebelum sekitar 31,4 s, sebagian besar panel gelap pada skala yang
dipilih. Setelah itu muncul jalur terang sekitar 0,8 kHz dan
kelipatannya, dengan perubahan posisi sebelum dan sesudah bagian
yang relatif mendatar sekitar 31,8--32,5 s. Hubungan mendekati
kelipatan tersebut mendukung keberadaan struktur harmonik, tetapi
tidak cukup untuk menetapkan fundamental atau mekanisme sumbernya.

Panel 4096 menampilkan garis lebih tipis; pada panel 1024 komponen
berdekatan lebih melebar. Pergeseran jalur memperlihatkan bahwa
komponen kuat tidak selalu berada pada frekuensi yang tetap.
Secara matematis, FFT global merangkum seluruh perubahan itu,
sedangkan STFT melokalisasinya melalui jendela. Kedua hasil berasal
dari sampel dekode yang sama, sehingga detail gambar tidak boleh
diartikan sebagai informasi audio baru yang muncul akibat konversi
MP3 ke WAV. Cuplikan kedua kanal tersedia dalam notebook untuk
pemeriksaan pendengaran oleh pembaca.

\subsection{Keterbatasan Penelitian}
FFT merangkum kontribusi frekuensi seluruh pengamatan; STFT
memperlihatkan lokalisasinya. Fourier kompleks lengkap tetap dapat
diinversikan. Yang tidak langsung memperlihatkan waktu adalah
grafik magnitudo atau dayanya, bukan berarti transformasi kompleks
menghilangkan semua informasi waktu. Spektrogram juga tidak
menampilkan fase STFT secara langsung.

Batas utama penelitian adalah pergantian preset dan mesin bunyi.
Partitur, velocity, register dan waktu aktif pun berbeda antarkelompok,
sehingga pengaruh jenis instrumen tidak diisolasi. Perataan stereo
mengurangi informasi spasial; RMS bukan tekanan suara; puncak
terbesar bukan identifikasi fundamental otomatis. Jendela tetap
tidak dapat menajamkan waktu dan frekuensi secara arbitrer sekaligus.

Cuplikan 29--33 s untuk lima instrumen dan campuran selaras, serta
cuplikan DRUMS dan VOCALS menurut waktu berkas sumber, tersedia
di folder \texttt{hasil/audio/cuplikan} serta pemutar notebook. Pembaca dapat
membandingkan audio dengan grafik. Uraian membedakan fakta komputasi,
interpretasi teori dan analogi bunyi; tidak mengaku melakukan uji
dengar subjektif.
